% =====================================================================
%  Technical appendix: full proofs of the deterministic edge certificate.
%  \input by main.tex after the bibliography. Restated results are kept
%  UNNUMBERED (bold) so the body's Prop.~1 / Thm.~1--3 numbering is intact.
% =====================================================================

\section{Proofs of the Edge Certificate}
\label{app:proofs}

Throughout, an explanation is the set $\Ek=\{e_1,\dots,e_k\}$ of the $k$
highest-scoring \emph{undirected} edges of the clean graph, ordered
$s(e_1)\ge\cdots\ge s(e_M)$. An adversary applies a perturbation
$\pi=(\mathcal{A},\mathcal{D})$: a set $\mathcal{A}$ of added edges and a set
$\mathcal{D}$ of deleted edges, with total budget
$|\mathcal{A}|+|\mathcal{D}|\le B$. Let $G'$ be the perturbed graph and
$\Ek'$ its top-$k$ explanation. All scores are taken from the Causal Discovery
Core, $s(u,v)=\sigma(\mathrm{MLP}[\,h_u,h_v,|h_u-h_v|,\cos(x_u,x_v)\,])$ with
$h_u=\mathrm{MLP}(x_u)$.

\paragraph{Undirected aggregation.}
The core scores directed edges; \aeth{} stores each physical edge as the pair
$(u,v),(v,u)$ and the unit of perturbation is the \emph{physical} (undirected)
edge. We assign each physical edge the mean of its two directed scores,
$s(\{u,v\})=\tfrac12\big(s(u,v)+s(v,u)\big)$. Because each directed score is a
function of its own two endpoints' features only, $s(\{u,v\})$ is a function of
$\{x_u,x_v\}$ only; adding or deleting any \emph{other} physical edge leaves it
unchanged. This is the sole property the proofs use.

\subsection*{Proposition~\ref{prop:local} (Edge-locality), restated}
\emph{For any perturbation that adds or deletes edges other than $e$, the score
$s(e)$ is unchanged.}

\begin{proof}
$s(e)$ depends only on the features of $e$'s two endpoints (above). Edge
additions and deletions do not alter node features, and the node set is fixed
(both threat models in \xgn{}/\pgn{} perturb $E$, not $X$ or $V$). Hence any
$\pi$ that does not itself add or delete $e$ leaves the features of $e$'s
endpoints---and therefore $s(e)$---invariant.
\end{proof}

A first corollary, used repeatedly: \emph{deletions of edges other than $e$
leave both $s(e)$ and the entire relative ranking of all surviving edges
unchanged}, since every surviving score is individually invariant and ranking
depends only on the multiset of scores.

\subsection*{Theorem~\ref{thm:overlap} (Certified overlap), restated}
\emph{For any $\pi$ with $|\mathcal{A}|+|\mathcal{D}|\le B$,
$\;|\Ek\cap\Ek'|\ge k-B$.}

\begin{proof}
We bound the number of clean explanatory edges that can leave the top-$k$. An
edge $e_i\in\Ek$ can fail to lie in $\Ek'$ in exactly two ways.

\emph{(i) Direct deletion.} If $e_i\in\mathcal{D}$, it is absent from $G'$ and
hence from $\Ek'$. Each deletion removes at most one member of $\Ek$, so
deletions account for at most $|\mathcal{D}|$ losses.

\emph{(ii) Eviction by an addition.} Suppose $e_i\notin\mathcal{D}$. By
Prop.~\ref{prop:local} its score $s(e_i)$ and the scores of all other surviving
clean edges are unchanged, so among the surviving clean edges the ranking is
identical to the clean ranking. The only edges that can outrank $e_i$ in $G'$
beyond those that already did are the added edges $\mathcal{A}$. Each added edge
occupies exactly one position in the sorted order of $G'$; thus the $a$ added
edges collectively push the clean edges down by at most $a$ positions in total,
and---because each added edge sits in a single slot---can displace at most
$|\mathcal{A}|$ clean edges out of the top-$k$. Formally, the number of clean
survivors with final rank $>k$ is at most the number of added edges with score
exceeding the clean $k$-th score, which is at most $|\mathcal{A}|$.

Combining, the number of members of $\Ek$ absent from $\Ek'$ is at most
$|\mathcal{D}|+|\mathcal{A}|\le B$. Therefore
$|\Ek\cap\Ek'|\ge k-B$.
\end{proof}

\subsection*{Theorem~\ref{thm:radius} (Per-edge certified radius), restated}
\emph{The rank-$r$ edge $e_r\in\Ek$ (with $r\le k$) remains in the top-$k$
explanation under any $\pi$ with $|\mathcal{A}|+|\mathcal{D}|\le R(e_r)=k-r$ that
does not delete $e_r$.}

\begin{proof}
Assume $e_r\notin\mathcal{D}$. By Prop.~\ref{prop:local}, $s(e_r)$ is unchanged
and every surviving clean edge keeps its score, so in $G'$ the only edges that
can be ranked above $e_r$ are (a) the $r-1$ clean edges already above it and (b)
those added edges whose score exceeds $s(e_r)$. Deletions of other edges only
remove competitors, weakly improving $e_r$'s rank. Writing $a\le|\mathcal{A}|$
for the number of added edges scoring above $s(e_r)$, the rank of $e_r$ in $G'$
is at most $r+a$. Hence $e_r$ stays in the top-$k$ whenever $r+a\le k$, i.e.\
whenever $a\le k-r$. Since $a\le|\mathcal{A}|\le B$, a budget $B\le k-r$
suffices, giving the certified radius $R(e_r)=k-r$.
\end{proof}

\noindent\textbf{Remark (consistency with Thm.~\ref{thm:overlap}).} Summing the
per-edge radii recovers the overlap bound: the $j$-th-from-bottom member of
$\Ek$ has rank $r=k-j+1$ and radius $j-1$, so it can be evicted only once the
budget exceeds $j-1$; at budget $B$ at most the bottom $B$ members are at risk,
again yielding $|\Ek\cap\Ek'|\ge k-B$.

\subsection*{Theorem~\ref{thm:hijack} (Transductive hijack impossibility), restated}
\emph{In the transductive setting, let the structural-prior allowlist be active
with cap $c<1$, so every edge outside the training adjacency has score at most
$c$. If every $e_i\in\Ek$ has clean score $s(e_i)>c$, then no inserted edge can
enter $\Ek'$, for any budget $B$; the explanation can lose members only by
direct deletion.}

\begin{proof}
Every added edge is, by construction, absent from the training adjacency, so the
allowlist multiplier caps its score at $c$. By hypothesis $s(e_i)>c$ for all
$e_i\in\Ek$, hence every added edge scores strictly below every clean
explanatory edge. In the proof of Thm.~\ref{thm:radius}, the count $a$ of added
edges outranking a member $e_i$ is therefore $0$ regardless of $|\mathcal{A}|$:
the certified radius against additions is $R(e_i)=\infty$. Consequently no added
edge ever occupies a top-$k$ slot, and the only way to reduce $|\Ek\cap\Ek'|$ is
to delete clean members directly. The attacker-\emph{insertion} (hijack) rate is
thus exactly $0$ for unbounded $B$.
\end{proof}

This guarantee has no analogue for voting-based certificates: a smoothing/voting
explainer can always be hijacked with enough budget because an inserted edge
participates in (and can flip the vote of) many sub-graphs. The empirical check
agrees---$18/18$ top-$k$ edges survive $B=100$ insertions.

\subsection*{Tightness}
The overlap bound is achieved with equality, so no sound certificate of this
form can be larger. Fix $B\le k$. Construct an adversary that adds $B$ edges
whose endpoint features yield scores above $s(e_k)$ (e.g.\ duplicating the
endpoint profile of $e_1$). By the eviction argument these displace exactly the
$B$ lowest-ranked members of $\Ek$, giving $|\Ek\cap\Ek'|=k-B$. Hence the bound
of Thm.~\ref{thm:overlap} is tight, the certified--empirical gap is $0$, and the
greedy adversary of Fig.~\ref{fig:tight} realizes it---in contrast to randomized
smoothing, whose lower bound is strictly below the worst case it admits.

\subsection*{Soundness gate}
Because the certificate is deterministic, soundness is checkable exactly. For
each trained model and dataset we run the worst-case adversary of the tightness
construction at every budget $B$ and confirm the \emph{realized} overlap never
falls below $k-B$ (function \texttt{soundness\_check}). This gate must pass
before any large certification sweep; an unsound certificate would otherwise
invalidate every downstream number. All reported configurations pass with the
gap closing exactly at the bound.
